\clearpage
\fancyhead[L]{\small GRÖSSTE SONNENHÖHE.}
\hypertarget{native-465-sunheight}{}
\begin{center}{\large GRÖSSTE SONNENHÖHE.}\end{center}

\begin{center}
\rule{0.55\linewidth}{0.4pt}\\[2pt]
Astronomische Hülfstafeln herausg. von \textsc{Schumacher} für 1827 Seite 92.\\[2pt]
\rule{0.55\linewidth}{0.4pt}
\end{center}

Bekanntlich erreicht die Sonne wegen der Änderung ihrer Deklination ihre größte Höhe außerhalb des Meridians, und zwar in den aufsteigenden Zeichen \emph{nach}, in den absteigenden Zeichen \emph{vor} ihrem Durchgang durch den Meridian. Der Zeitpunkt dieser größten Höhe ist bemerkenswert, weil er nach einer einfallsreichen Bemerkung von Hofrat \textsc{Gauss} bei Beobachtungen
\SourceLeaf{466}meridiannaher Sonnenhöhen dem Berechner die Mühe erspart, die Änderung der Deklination während der Beobachtungszeit zu berücksichtigen. Die Stundenwinkel dienen als Argumente der Korrekturen, mit denen man die zu anderen Zeiten als am Mittag gemessenen Höhen auf die Mittagshöhe reduziert. Man braucht die Stundenwinkel lediglich vom Zeitpunkt der größten Sonnenhöhe aus zu rechnen statt wie gewöhnlich vom wahren Mittag aus. Mit dieser einfachen Verschiebung des Nullpunkts der Stundenwinkel ist der Einfluss der Deklinationsänderung für die gesamte Beobachtungsreihe berücksichtigt. --- Die folgende Tafel gibt für jeden zehnten Tag und für alle Polhöhen von Gibraltar bis Petersburg (von $36^\circ$ bis $60^\circ$) an, was man zum Zeitpunkt des wahren Mittags hinzuzählen oder von ihm abziehen muss, um den Zeitpunkt der größten Sonnenhöhe zu erhalten. \textsc{Schumacher.}
